\scanpage{340}
% @ top centre
\mspage{64}

% @ top, item 1 and interline
1) La théorie des \uline{catégories de chemins} $\underline{\mathrm{Ch}}(X)$ ($\underline{\overline{\mathrm{VII}}}$), \inserted{s'appuyant sur celle des \uline{morphismes lisses}, \uline{propres}\unsure{,} \uline{parfaits}.}

% @ top, item 2
2) Le lemme de K.S. Brown (3.2.1 \unsure{pg|pag|page} 20)

% @ upper third, item 3
3) Pour l'implication $B\Rightarrow A$ i.e. $\mathrm{Loc}(7\mathrm b)\Rightarrow \mathrm{Loc}\,7$, (sous l'hypothèse $W(1\ \mathrm{ter},\deleted{\unsure{8\mathrm{bis}}}3\ \mathrm{bis})$), j'utilise la théorie du prédérivateur $\mathrm{HOT}_W$ (VI), et plus particulièrement celle des foncteurs $f_{\unsure{!|.}}$ (définis sous la seule hypothèse que \deleted{$W$ satisfasse} \deleted{\unsure{$\mathrm{Loc}$}}$W(3\ \mathrm{bis})$, renforçant la condition automatique $W(3)$ (faisant partie de la définition d'un loc fond.)

% @ centre and lower middle, Brown and open question
Il est clair que je ne serais pas arrivé à y voir clair, sans le travail fondamental de K.S. Brown. Il reste pourtant encore des questions ouvertes. Ainsi, si on ne suppose pas que $W$ satisfasse à $\mathrm{Loc}(8\ \mathrm{bis})$, mais on \unsure{borne|se borne|tient} à $\mathrm{Loc}(3\ \mathrm{bis})$ (et je n'ai pas eu l'occasion encore de travailler avec un localiseur qui n'y satisfasse \ldots), est-il vrai néanmoins que l'on a \inserted{encore} l'implication

% @ lower quarter, question's implication
\[
 D\Longrightarrow A,
\]
% @ foot, question continues
de façon que les conditions $A,B,C,D$ sont équivalentes (et en particulier $\mathrm{Loc}(7\mathrm c)$ implique déjà $\mathrm{Loc}(7)$ dans \deleted{\unsure{toute}} sa force) ? Il est clair dans tout cela que
% archive sheet number: 339
